\section{Conclusion}

Le modèle de Bergeron associé au modèle compagnon de la règle des trapèzes retrouve les valeurs
du livre ($i_2 = \ideuxC$~A à \SI{0.7}{ms}) et le régime permanent calculé par phaseurs. L'onde
atteint l'extrémité $m$ après $\tau = \tauL~\unit{\micro s}$. La tension y double presque à
l'extrémité ouverte (\rapMega{}\,$V_m$ avec \SI{e6}{\ohm}, mais seulement \rapNeuf{}\,$V_m$ avec
999~\unit{\ohm}) et s'effondre en court-circuit. Avec $\Delta t = \SI{0.1}{ms}$, l'interpolation
du retard amortit les oscillations libres ; un pas $\Delta t = \tau/N$ éviterait cet
amortissement artificiel. Pour le choc de foudre, la ligne de 100~mi n'a pas le temps de réagir
en \SI{250}{\micro s}, et il faut la raccourcir à 10~mi pour voir le doublement de tension
(\gVRL~kV).
